\begin{exercise}{119}
Prove Lemma~3.2.7

\begin{proof}
Let $p\colon X\to S$ be an inner fibration of simplicial sets.
Let $\sigma\colon\Delta^2\to X$ be the given triangle, with
$f=d_2\sigma$, $g=d_0\sigma$, $h=d_1\sigma$, and $f$ $p$-cocartesian.
Pull $p$ back along $p\sigma\colon\Delta^2\to S$, obtaining
\[
 q\colon X_\sigma=X\times_S\Delta^2\longrightarrow\Delta^2.
\]
The map $q$ is an inner fibration, so both its source and target are
$\infty$-categories.
The canonical lift $(\sigma,\mathrm{id}_{\Delta^2})$ has edges
$\widetilde f,\widetilde g,\widetilde h$ projecting to $f,g,h$.
By the proof of Exercise~\ref{ex:117}, $\widetilde f$ is $q$-cocartesian,
and so is $\widetilde g$ or $\widetilde h$ whenever the corresponding
original edge is $p$-cocartesian.

The dual of Lemma~3.1.7, applied over the $\infty$-category $\Delta^2$,
gives that $\widetilde g$ is $q$-cocartesian if and only if
$\widetilde h$ is $q$-cocartesian.
Thus either hypothesis of Lemma~3.2.7 makes the other lifted edge
$q$-cocartesian.
Exercise~\ref{ex:118}, applied to this pullback, then makes its image
locally $p$-cocartesian in $X$.
This proves both asserted implications.
\end{proof}
\end{exercise}
